\begin{abstract}
We introduce \emph{tactile analogies}, an example-based synthesis paradigm for tactile videos.
Given a small set of reference touches recorded on an object, we predict the tactile signal that a
visuo-tactile sensor would record at a new, previously untouched location. We represent a touch as
a \emph{tactile normal video}, which records how the gel surface of the sensor deforms over the
course of a press, and we bridge the reference and query locations through surface normal maps
rendered at the two sensor poses. Our coarse-to-fine pipeline (i) compares 2D foundation model features to choose the most promising
reference touch, (ii) aligns the two renderings with a pre-trained local feature matcher and warps
the reference video into the query frame, and (iii) edits the warped video into the final
prediction with a lightweight refinement network.
As the task is new, we build a benchmark on a large database of 3D objects, on which our method
outperforms baselines adapted from prior work in tactile prediction while running in 0.6\,s per
prediction, making it suitable for interactive applications in AR/VR and robotics. We further
outline how the predicted normals could be integrated into heightmaps to estimate 3D surface
geometry, and passed through an optical tactile simulator to approximate visuo-tactile measurements
at locations that were never touched.
\end{abstract}
